%% ====================================================================
\section{Introducción}
\label{sec:intro}
%% ====================================================================

La operación de un almacén automatizado obliga a coordinar flotas heterogéneas frente
a cargas que demandan distinto número de manipuladores. No todas pesan lo mismo. En un
escenario con quince AGV y cinco tipos de carga activos en simultáneo (la
Figura~\ref{fig:problema} ilustra un subcaso con $N{=}6$ y $K{=}3$), un solo robot basta
para cada una de tres cajas ligeras, una paleta exige dos y una estructura de acero,
cuatro. Las cargas aparecen y se completan sin pausa. La flota debe reasignarse y formar
coaliciones que satisfagan la cardinalidad mínima de cada tarea. El
resultado es una asignación distribuida con restricciones de cardinalidad que deben
cumplirse en tiempo real, sin coordinador central y con comunicación limitada al entorno
inmediato de cada robot. Difiere así de la planificación de rutas y del seguimiento de
formación fija.

Frente al transporte cooperativo clásico, donde varios robots empujan juntos una sola
carga \citep{ebel2024cooperative, rosenfelder2024force}, este problema invierte la
dificultad. Allí la composición del grupo está fijada de antemano; el reto es controlar
el conjunto. Aquí esa composición es la variable de diseño. La flota debe resolver
localmente cuántos robots asigna a cada tarea, y la decisión solo vale si produce
coaliciones factibles dentro del horizonte de operación. Los métodos de
asignación distribuida por subasta \citep{choi2009consensus} o por formación
de coaliciones \citep{dutta2021hedonic, shan2024distributed} abordan la
coordinación desde una perspectiva discreta o negociada; en la literatura revisada,
estos enfoques no se formulan como una dinámica poblacional continua acoplada a
navegación reactiva para inducir coaliciones de cardinalidad mínima. El aprendizaje
por refuerzo \citep{bezerra2025learning} produce políticas
efectivas pero dependientes de la distribución de entrenamiento y sin garantías
de robustez ante configuraciones de tarea no vistas.

\subsection{Dinámicas poblacionales como ley de control distribuida}

Una línea de investigación alternativa, desarrollada principalmente por el grupo
de Quijano, propone emplear dinámicas evolutivas de juegos poblacionales
directamente como leyes de control distribuido \citep{quijano2017population}.
La proporción de robots asignados a cada tarea se lee como el perfil de estrategias de
una población, y una dinámica de revisión actualiza ese perfil según las diferencias de
\emph{payoff} (función de utilidad que recibe cada robot por elegir una tarea). El
coordinador central desaparece. Cada agente actúa solo con información local, y el
sistema converge al equilibrio Nash bajo condiciones de juego potencial.

Como muestran \citet{barreiro2017distributed} y \citet{martinez2020formation}, las
dinámicas distributivas del replicador permiten coordinar formaciones de robots móviles
con comunicación limitada en grafo, extensión que \citet{barreiro2021uav} lleva a
formaciones de UAVs con datos reales. Más recientemente, \citet{dinh2026bsplines} proponen
mecanismos basados en B-splines para diseñar funciones de payoff que satisfagan
restricciones de formación específicas. En esa línea, sin embargo, suele suponerse que
el tamaño del grupo que realiza cada tarea es fijo y conocido de antemano. El requisito
de cardinalidad mínima dinámica no recibe tratamiento explícito.

\subsection{El papel de la política de revisión: replicador y Smith}
\label{subsec:porquesmith}

Manteniendo fija la arquitectura de estimación, una de las preguntas que la escalera
comparativa del trabajo formaliza es qué aporta la \emph{política de revisión} de una
dinámica poblacional. Los tratamientos T3 (replicador distribuido)
y T4 (dinámicas de Smith) comparten el mismo payoff y el mismo estimador por
consenso, y difieren únicamente en esta ley; comparar ambos aísla su efecto. Ambas
parten del mismo marco. El estado de la población es el vector de fracciones
$\mathbf{x} = (x_k)$, donde $x_k$ es la proporción de robots asignados a la tarea $k$.
Cada robot estima $x_k$ por consenso promedio distribuido y obtiene así
$\hat{x}_{ik}$. Hasta aquí, Smith y DRD son idénticos: ambos requieren esta
\emph{primera capa de consenso}.

La diferencia aparece en lo que cada dinámica requiere más allá de esa primera capa.
Las dinámicas del replicador (DRD) comparan el payoff de cada estrategia con la
\emph{media poblacional} $\bar{f}(\mathbf{x}) = \sum_k x_k f_k(\mathbf{x})$:
esta cantidad es global porque depende de la distribución completa $\mathbf{x}$,
no solo del estado local de un agente. Estimarla distribuida\-mente exige una
\emph{segunda capa de consenso} sobre los payoffs individuales $f_k(\hat{x}_i)$.
Cuando esta segunda capa opera sobre un grafo irregular que cambia con el movimiento
de los robots, el operador de consenso actúa como un filtro pasa-bajos espectral
\citep{sandryhaila2013dsp}: atenúa las componentes de alta frecuencia del vector
de payoffs sobre la estructura del Laplaciano de $G(t)$ e introduce un sesgo
dependiente de la topología $\delta(G, \beta)$ que desplaza el equilibrio efectivo
respecto al equilibrio Nash verdadero.

A diferencia del replicador, las dinámicas de Smith evitan esta segunda capa porque la
actualización se basa en comparaciones bilaterales $\pos{f_{ik} - f_{ij}}$ calculadas localmente con los
payoffs propios: ningún agente necesita estimar la media poblacional global. La
condición de punto fijo $\pos{f_{ik} - f_{ij}} = 0$ se satisface exactamente con
información local derivada de $\hat{x}_{ik}$. Además, las estrategias extintas
($p_{ik} = 0$) pueden recuperarse, ya que la tasa de migración es independiente
de la masa actual en cada estrategia. La línea base T3 (tratamiento de replicador
distribuido) cuantifica empíricamente el efecto de este sesgo bajo las mismas
topologías de comunicación que el método propuesto.

\subsection{La contribución de este trabajo}

La pieza central del trabajo es una función de payoff con componente sigmoide decreciente
que representa requisitos mínimos de cardinalidad mediante cotas inferiores, en contraste
con las cotas superiores de \citet{martinezpiazuelo2022automatica}. En el caso base con
$\Psi\equiv1$ preserva una estructura de juego de congestión ponderado. Al activar
$\Psi(d)$, el sistema implementado se vuelve una perturbación espacial que acopla
asignación y movimiento. Sobre ese payoff se adaptan y comparan el replicador distribuido
y las dinámicas de Smith bajo una misma arquitectura de estimación por consenso dinámico
de media; así se aísla el efecto de la política de revisión. Una tasa de revisión
espacialmente modulada suprime, además, las reasignaciones espurias durante la
aproximación a las cargas y en régimen estacionario. Queda un último paso: medir el coste
de la descentralización frente a una referencia centralizada replanificada y la
degradación al reducir el rango de comunicación. Estos cuatro elementos definen una
arquitectura donde la estimación de ocupación, la revisión estratégica y la navegación
reactiva se actualizan en un mismo bucle temporal, sin coordinador central de asignación.

La validación se realiza mediante simulación cuantitativa en Python con seis
escenarios y cinco tratamientos comparativos dispuestos como escalera incremental
(T1 heurística voraz local, T2 DAC con regla voraz, T3 replicador distribuido
modificado, T4 dinámicas de Smith con revisión modulada (método propuesto) y
T5 referencia centralizada replanificada), más una validación cualitativa y
cinemática en CoppeliaSim.

\subsection{Estructura del documento}

La sección~\ref{sec:planteamiento} describe el problema de asignación distribuida
con restricciones de cardinalidad y la brecha que lo motiva. La
sección~\ref{sec:hipotesis} formula la pregunta de investigación y las hipótesis.
Los objetivos se especifican en la sección~\ref{sec:objetivos}, y el alcance y
limitaciones en la sección~\ref{sec:alcance}. La sección~\ref{sec:metodologia},
núcleo de esta entrega, desarrolla las tres dinámicas acopladas, la
función de payoff, el análisis de convergencia y el diseño experimental completo.
Las secciones \ref{sec:marco}--\ref{sec:conclusiones} están reservadas para el
marco teórico, resultados y conclusiones de la entrega final.

\subsection{Estado actual del TFM}
\label{subsec:estado}

A la fecha de esta entrega (mayo de 2026) se han completado las siguientes etapas:

\begin{itemize}
  \item \textbf{Revisión bibliográfica} (21 referencias en 6 áreas temáticas:
    dinámicas poblacionales, consenso distribuido, formación de coaliciones,
    transporte cooperativo, juegos potenciales y procesado de señales en grafos).
  \item \textbf{Formulación matemática completa}: modelo de robots, tareas y
    grafo de comunicación; función de payoff multiplicativa con sigmoide
    decreciente; tres dinámicas acopladas (consenso, Smith en tiempo discreto,
    gradiente reactivo); análisis formal de convergencia para el caso simétrico
    (Teorema~\ref{prop:convergencia} con demostración y cálculo de Hessiana).
  \item \textbf{Diseño experimental}: seis escenarios, cinco tratamientos
    comparativos (incluido el replicador distribuido como línea base clave),
    once métricas definidas y criterios de aceptación conjuntivos.
  \item \textbf{Pseudocódigo formal} del bucle de control distribuido
    (Algoritmo~\ref{alg:loop}).
\end{itemize}

En curso: implementación en Python de las tres dinámicas acopladas, los
tratamientos comparativos y la campaña experimental. El plan de trabajo
completo se detalla en la Sección~\ref{sec:cronograma}.

\subsection{Planteamiento del problema y justificación}
\label{sec:planteamiento}

\subsubsection{Descripción formal del problema}

El Anexo~I de este trabajo mencionaba ``cargas heterogéneas'' (objetos que difieren
en peso, fragilidad y tamaño). Aquí esa heterogeneidad se formaliza
como un requisito mínimo de cardinalidad $n_k$: el número de manipuladores
que una carga determinada exige para ser transportada con seguridad. La abstracción es
intencional: atributos como la masa o la fragilidad se condensan en un valor escalar
$n_k$ que el sistema de asignación consume; derivar $n_k$ a partir de un modelo de contacto explícito (cierre de
fuerza, estabilidad del agarre) es una capa complementaria documentada en
\citet{ebel2024cooperative} y \citet{rosenfelder2024force} que queda fuera del
alcance de este TFM (véase \S\,\ref{sec:alcance}). El problema central y el objeto de
estudio del Anexo~I se mantienen; este informe concreta la terminología (formación de
coaliciones) y la metodología adoptada.

Con esta abstracción, el sistema se define así. Se considera
una flota de $N$ AGV que operan en un entorno 2D. El conjunto de robots
se indexa $i \in \{1, \ldots, N\}$ y el conjunto de tareas activas en un instante
dado se indexa $k \in \{1, \ldots, K\}$. Cada tarea $k$ está ubicada en una
posición conocida $\mathbf{l}_k \in \RR^2$ y exige que al menos $n_k \in \NN$
robots lleguen a ella para que la carga pueda ser manipulada. La tarea $k=0$ representa
el estado de inactividad: un robot inactivo (\textit{idle}) no contribuye a ninguna tarea.

Cada robot mantiene un vector de estrategias $\mathbf{p}_i(t) \in \simplex^K$,
donde $\simplex^K = \{p \in \RR^{K+1} : \sum_{k=0}^K p_k = 1,\, p_k \geq 0\}$
es el símplex de probabilidad. El componente $p_{ik}(t)$ representa la propensión
del robot $i$ a seleccionar la tarea $k$. La estrategia efectiva del robot es
$s_i(t) = \argmax_k p_{ik}(t)$. El objetivo del sistema es que, para cada tarea $k$,
la coalición $\mathcal{C}_k(t) = \{i : s_i(t) = k, \norm{\mathbf{q}_i - \mathbf{l}_k}
< r_{\text{det}}\}$ satisfaga $|\mathcal{C}_k(t)| \geq n_k$ en tiempo finito, y
que esta condición se mantenga robustamente frente a cambios en el conjunto de
tareas activas.

\begin{figure}[htbp]
\centering
\begin{tikzpicture}[x=1cm, y=1cm,
  agv/.style={circle, draw=VIUOrange, fill=VIUOrange!20,
    minimum size=0.62cm, font=\scriptsize\bfseries, inner sep=0pt},
  pickup/.style={rectangle, draw=VIUDark, fill=VIUDark!8,
    minimum width=0.88cm, minimum height=0.88cm, font=\tiny, align=center, inner sep=2pt},
  comm/.style={draw=VIUGray!70, thin, dashed},
  motion/.style={-{Stealth[length=4pt]}, draw=VIUOrange!75, thick, shorten >=3pt, shorten <=3pt}
]
%-- Planta de bodega
\draw[VIUDark!40, thick, rounded corners=5pt] (0,0) rectangle (11.5,7.8);
\node[font=\tiny, VIUGray, anchor=north west] at (0.15,7.65)
  {Almacén automatizado ($20\,\mathrm{m}\times 14\,\mathrm{m}$)};

%-- Tareas: posiciones de recogida l_k
\node[pickup] (L1) at (2.0,6.5) {$k{=}1$\\$n_1{=}1$};
\node[pickup] (L2) at (5.5,3.2) {$k{=}2$\\$n_2{=}3$};
\node[pickup] (L3) at (9.5,6.0) {$k{=}3$\\$n_3{=}2$};
\node[font=\tiny, VIUDark, above=1pt of L1] {$\mathbf{l}_1$};
\node[font=\tiny, VIUDark, above=1pt of L2] {$\mathbf{l}_2$};
\node[font=\tiny, VIUDark, above=1pt of L3] {$\mathbf{l}_3$};

%-- Destinos d_k (dibujados como rombos manuales)
\def\rdia{0.30}
\foreach \cx/\cy/\lbl in {1.5/1.0/$\mathbf{d}_1$, 8.5/0.8/$\mathbf{d}_2$, 10.8/3.8/$\mathbf{d}_3$}{%
  \draw[VIUDeep, thick]
    (\cx,\cy+\rdia) -- (\cx+\rdia,\cy) -- (\cx,\cy-\rdia) -- (\cx-\rdia,\cy) -- cycle;
  \node[font=\tiny, VIUDeep] at (\cx,\cy) {\lbl};
}

%-- Robots
\node[agv] (R1) at (2.0,5.0)  {1};
\node[agv] (R2) at (4.4,2.7)  {2};
\node[agv] (R3) at (6.5,2.7)  {3};
\node[agv] (R4) at (5.6,4.2)  {4};
\node[agv] (R5) at (8.5,5.0)  {5};
\node[agv] (R6) at (9.8,4.8)  {6};

%-- Coalicion C_2 = {2,3,4}: recuadro resaltado
\draw[VIUOrange, thick, dashed, rounded corners=6pt, fill=VIUOrange!5]
  (3.7,2.0) rectangle (7.2,5.0);
\node[font=\tiny, VIUOrange, anchor=south west] at (3.7,5.0)
  {$\mathcal{C}_2$: $|\mathcal{C}_2|=3\geq n_2$ $\checkmark$};

%-- Flechas de transporte: coalicion 2 se mueve hacia d_2
\draw[{Stealth[length=4pt]}-{Stealth[length=4pt]}, VIUOrange!60, thick, dashed]
  ($(L2.south)+(0,-0.1)$) -- (8.5,0.8)
  node[midway, below, font=\tiny, VIUGray] {transporte $\rightarrow \mathbf{d}_2$};

%-- Links de comunicacion dentro de la coalicion
\draw[comm] (R2) -- (R3);
\draw[comm] (R2) -- (R4);
\draw[comm] (R3) -- (R4);

%-- Robots 5 y 6: convergiendo hacia L3
\draw[comm] (R5) -- (R6);
\draw[motion] (R5) -- ($(L3.south west)+(0.05,0.05)$);
\draw[motion] (R6) -- ($(L3.east)+(-0.05,0)$);

%-- Robot 1: solo, hacia L1
\draw[motion] (R1) -- ($(L1.south)+(0,0.05)$);

%-- Circulo de rango R (ejemplo alrededor de R5)
\draw[VIUGray, dotted, thin] (R5) circle (1.4cm);
\node[font=\tiny, VIUGray, anchor=north] at ($(R5)+(0,-1.4)$) {rango $R$};

%-- Leyenda
\begin{scope}[shift={(0.1,0.05)}]
  \node[agv, minimum size=0.42cm] at (0.38,0.38) {};
  \node[font=\tiny, anchor=west] at (0.65,0.38) {AGV $i$};
  \node[pickup, minimum width=0.42cm, minimum height=0.42cm] at (2.5,0.38) {};
  \node[font=\tiny, anchor=west] at (2.8,0.38) {Recogida $\mathbf{l}_k$};
  \draw[VIUDeep, thick]
    (5.2,0.62) -- (5.45,0.38) -- (5.2,0.14) -- (4.95,0.38) -- cycle;
  \node[font=\tiny, anchor=west] at (5.55,0.38) {Entrega $\mathbf{d}_k$};
  \draw[comm] (7.3,0.38) -- (8.0,0.38);
  \node[font=\tiny, anchor=west] at (8.05,0.38) {Enlace com.};
\end{scope}
\end{tikzpicture}
\caption[Escenario de asignación distribuida con coaliciones]%
{Escenario del problema de asignación distribuida con coaliciones. Seis AGV
($N=6$) en un almacén automatizado con tres tareas activas ($K=3$). La coalición $\mathcal{C}_2$
(robots~2, 3, 4) ya satisface $|\mathcal{C}_2|=3=n_2$ y transita hacia el destino
$\mathbf{d}_2$. Los robots~5 y~6 convergen hacia la tarea~3 ($n_3=2$). El robot~1 atiende
solo la tarea~1 ($n_1=1$). El círculo punteado indica el rango de comunicación~$R$.}
\label{fig:problema}
\end{figure}


Dos dificultades distinguen este problema de la asignación estática de tareas. La
restricción $|\mathcal{C}_k| \geq n_k$ es una \emph{cota inferior} dinámica: si tres
robots abandonan una tarea que requiere cuatro, el sistema debe reclutar uno nuevo sin
coordinación central. A ello se suma que, cuando las tareas aparecen, desaparecen o
cambian su prioridad durante la operación, las coaliciones deben reconstituirse sin
reinicializar el sistema completo.

\subsubsection{Límites de las soluciones existentes}

Los métodos de asignación descentralizada basados en subastas
\citep{choi2009consensus} y en formación de coaliciones hedo\-nistas
\citep{dutta2021hedonic, zhang2024coalition} coordinan la asignación mediante
negociación iterativa entre agentes, pero producen asignaciones discretas que
no se actualizan continuamente ante cambios en el entorno. El aprendizaje por refuerzo
\citep{bezerra2025learning} genera políticas efectivas en distribución, pero exige
entrenamiento previo sobre escenarios representativos y carece de garantías formales de
convergencia ante configuraciones de tarea no vistas. Ambas familias comparten un mismo
límite. No proporcionan una ley de control diferencial que actualice sin interrupción
las estrategias de asignación en tiempo real.

En cambio, las dinámicas poblacionales \citep{barreiro2021uav, dinh2026bsplines,
hurtado2024potential} sí ofrecen leyes de control con garantías de convergencia
al equilibrio Nash, gracias a la estructura de juego potencial subyacente. Aun así,
en todos los trabajos de esta línea el tamaño del grupo que realiza
cada tarea es un parámetro del diseño fijado de antemano: los robots se asignan
en fracciones fijas, no en coaliciones cuyo tamaño emerge de la dinámica.
Cuando el tamaño de la coalición pasa a ser la variable de estado (porque cada carga
impone su propio requisito mínimo), las dinámicas deben reclutar o liberar
robots según las necesidades de cada carga; ninguna formulación con ese grado de
generalidad apareció en los trabajos analizados.

El trabajo de \citet{martinezpiazuelo2022automatica} introduce restricciones de
capacidad en juegos poblacionales, pero estas restricciones son \emph{cotas superiores}
sobre cuántos agentes pueden elegir una estrategia: impiden que una estrategia se
sature por encima de un umbral. El caso que aquí se aborda requiere cotas inferiores:
la tarea no puede ejecutarse si hay \emph{menos} de $n_k$ robots. El cambio de signo no
es cosmético. Los resultados de estabilidad de ese trabajo no se trasladan directamente
a este escenario.

En las referencias consultadas tampoco se identifica un enfoque que trate
explícitamente el caso $N < \sum_k n_k$ (menos robots que los requeridos en total):
en ese régimen la
flota no puede satisfacer todos los requisitos de cardinalidad simultáneamente y
el sistema debe priorizar. Este caso límite, relevante en escenarios con alta
demanda transitoria, queda fuera del alcance de este trabajo pero se identifica como
línea futura.

\subsubsection{La brecha específica}

La combinación de tres elementos (dinámicas de Smith como ley de control directa,
función de payoff con cota inferior de cardinalidad mediante sigmoide decreciente
y tasa de revisión espacialmente modulada) no aparece en el corpus
bibliográfico revisado. Cada elemento por separado tiene antecedentes:
las dinámicas de Smith en sistemas de control \citep{barreiro2017distributed},
las funciones sigmoide para incentivos de reclutamiento en antropología
e inteligencia de enjambre \citep{theraulaz1999stigmergy},
y las tasas de revisión moduladas por contexto en juegos evolutivos
\citep{sandholm2010population}. La integración de los tres, con el análisis
formal que ella requiere, constituye la contribución de este trabajo.

La afirmación del trabajo es acotada a propósito: bajo comunicación local y
cargas heterogéneas con requisitos mínimos de cardinalidad, ¿puede una arquitectura
distribuida e interpretable que integre en un único bucle la estimación de ocupación,
la revisión poblacional de estrategias y la navegación reactiva formar coaliciones
factibles con una pérdida de desempeño acotada frente a una referencia centralizada?
Responderla es el objetivo. Afirmar la superioridad del control distribuido sobre el
centralizado en todo régimen operativo queda fuera de su alcance.

